\documentclass[UTF8,a4paper,11pt]{ctexart}

\usepackage{amsmath,amssymb}
\usepackage{geometry}
\geometry{margin=2.5cm}
\usepackage{booktabs}
\usepackage{hyperref}
\usepackage{enumitem}

\title{Fourier Analysis: From Summation to Integration}
\author{}
\date{}

\begin{document}
\maketitle

The transition from a \textbf{summation} ($\sum$) to an \textbf{integral} ($\int$) happens when the number of terms becomes very large and the spacing between them becomes infinitesimally small --- a key idea in calculus and applied math (e.g., Fourier analysis, probability, physics).

The picture you shared has \textbf{``Summation Discrete''} at the bottom, pointing toward \textbf{``Continuous''} above $\int$ --- this visually captures the idea:

\bigskip\noindent\rule{\linewidth}{0.4pt}\bigskip

\section{Summation of discrete terms}

A finite or infinite sum over discrete indices $k$:

\[
\sum_{k=1}^{n} f(x_k) \, \Delta x
\]

If $\Delta x$ is the spacing between successive $x_k$ values,
then as $n \to \infty$ and $\Delta x \to 0$,
the \textbf{Riemann sum} becomes an integral:

\[
\sum_{k=1}^{n} f(x_k) \, \Delta x \quad \longrightarrow \quad \int_{a}^{b} f(x) \, dx
\]

\bigskip\noindent\rule{\linewidth}{0.4pt}\bigskip

\section{Physical intuition (from your notes)}

Your sheet has a sequence:

\[
A_1 \to A_2 \to A_3 \to \dots
\]
and
\[
A_1 + A_2 + A_3 + \dots
\]

If these $A_n$ represent amplitudes (like in Fourier series), then:

\begin{itemize}
  \item \textbf{Discrete Fourier series}:
  \[
  f(t) = \sum_{n=-\infty}^{\infty} c_n e^{in\omega_0 t}
  \]

  \item As the period $T \to \infty$, the fundamental frequency $\omega_0 = 2\pi/T \to 0$,
  discrete frequency steps merge into a continuum.
  Then:
  \[
  \sum_{n=-\infty}^{\infty} c_n e^{in\omega_0 t} \quad \longrightarrow \quad \int_{-\infty}^{\infty} F(\omega) e^{i\omega t} \, d\omega
  \]
\end{itemize}

This is exactly the \textbf{Fourier transform} derived from the Fourier series by turning the sum into an integral.

\bigskip\noindent\rule{\linewidth}{0.4pt}\bigskip

\section{Why do this?}

\begin{itemize}
  \item Integrals handle \textbf{continuous spectra} (e.g., a non-periodic signal).
  \item Sums handle \textbf{discrete spectra} (periodic signals).
  \item The equation $\sum \to \int$ unifies both cases in the limit.
\end{itemize}

\bigskip\noindent\rule{\linewidth}{0.4pt}\bigskip

So in your notation:
\[
\sum \text{(Discrete)} \xrightarrow{\text{limit of small spacing}} \int \text{(Continuous)}
\]

Thanks for sharing the second image --- it contains a lot of handwritten notes that blend \textbf{Fourier transforms}, \textbf{Laplace transforms}, and \textbf{damped oscillations}. Let me break it down piece by piece.

\bigskip\noindent\rule{\linewidth}{0.4pt}\bigskip

\section{Top-left: Limit \& Fourier Transform}

\noindent\textit{lim $t \to \infty$ then Fourier Transform}

This suggests: as time grows large, we often look at frequency content using the Fourier transform.

The Fourier transform given is:

\[
\int_{-\infty}^{\infty} F(t) e^{-j\omega t} \, dt
\]

This transforms a time-domain signal $F(t)$ into the frequency domain.

\bigskip\noindent\rule{\linewidth}{0.4pt}\bigskip

\section{Euler's formula (center-left)}

\[
e^{j\omega t} = \cos(\omega t) + j\sin(\omega t)
\]
\[
e^{-j\omega t} = \cos(\omega t) - j\sin(\omega t)
\]

These are used to represent oscillatory behavior in terms of complex exponentials.

\bigskip\noindent\rule{\linewidth}{0.4pt}\bigskip

\section{Damped oscillator equation (center)}

\[
m\ddot{x} + b\dot{x} + Kx = 0
\]

That's a \textbf{spring-mass-damper} system.

The solution is written:

\[
x(t) = e^{-jwt} (A \sin(\omega t + \phi))
\]

\emph{Wait} --- this has a small error:
For a damped oscillator, the correct form is usually:

\[
x(t) = e^{-\zeta \omega_n t} \big(A \cos(\omega_d t) + B \sin(\omega_d t)\big)
\]

Here they wrote $e^{-j\omega t}$ instead of a real exponential decay $e^{-\sigma t}$. That suggests they are either:

\begin{itemize}
  \item Confusing \textbf{oscillation} $e^{j\omega t}$ with \textbf{decay}, or
  \item Using a \textbf{complex frequency} $s = \sigma + j\omega$ such that $e^{-j\omega t}$ is just the oscillatory part, while the real decay is hidden in the amplitude $A$ (not ideal notation).
\end{itemize}

\bigskip\noindent\rule{\linewidth}{0.4pt}\bigskip

\section{Laplace transform (right side)}

Two equivalent forms written:

\[
\int_{0}^{\infty} f(t) e^{-st} \, dt
\]
and
\[
\int_{0}^{\infty} F(s) e^{-st} \, ds
\]

The second one is \textbf{incorrect} --- that's the inverse Laplace transform.
Actually, the forward Laplace is:

\[
F(s) = \int_{0}^{\infty} f(t) e^{-st} \, dt
\]

\bigskip\noindent\rule{\linewidth}{0.4pt}\bigskip

\section{S-plane diagram (bottom right)}

They draw the \textbf{complex s-plane}:

\begin{itemize}
  \item Horizontal axis: $\text{Re}$ (real part $= \sigma$) --- corresponds to \textbf{growth/decay}
  \item Vertical axis: $\text{Im}$ (imag part $= \omega$) --- corresponds to \textbf{oscillation frequency}
\end{itemize}

This is where Laplace transforms live:
$s = \sigma + j\omega$

\bigskip\noindent\rule{\linewidth}{0.4pt}\bigskip

\section{Bottom-left ``Simulation Discrete / Summation Continuous''}

This is hard to read fully, but likely refers to:

\begin{itemize}
  \item \textbf{Discrete simulation} = numerical approximation (e.g., discrete time steps)
  \item \textbf{Summation $\to$ Continuous} = replacing a sum with an integral (connecting back to your first image)
\end{itemize}

They seem to be showing the \textbf{two perspectives}:

\begin{itemize}
  \item Discrete: $\sum$
  \item Continuous: $\int$
\end{itemize}

\bigskip\noindent\rule{\linewidth}{0.4pt}\bigskip

\section{Overall meaning of the picture}

This is a \textbf{summary sheet} connecting:

\begin{enumerate}
  \item \textbf{Fourier transform} (oscillations only, $s = j\omega$, no decay)
  \item \textbf{Laplace transform} (oscillations + growth/decay, $s = \sigma + j\omega$)
  \item \textbf{Mechanical oscillator} as the physical example
  \item \textbf{Discrete vs continuous} summation/integration as a unifying theme
\end{enumerate}

It's messy but shows the thought process of someone linking:
\begin{itemize}
  \item Time domain $\leftrightarrow$ frequency domain
  \item Summation $\leftrightarrow$ integration
  \item Real exponentials $\leftrightarrow$ complex exponentials
\end{itemize}

Would you like me to redraw this cleanly with the corrected equations?

\end{document}
